% Generated draft: requires mathematical and source review.
\documentclass[11pt]{article}
\usepackage[margin=25mm]{geometry}
\usepackage{fontspec}
\setmainfont{Latin Modern Roman}
\usepackage{amsmath,amssymb,mathtools,unicode-math}
\setmathfont{Latin Modern Math}
\usepackage{enumitem,needspace,xurl,hyperref}
\setlength{\parindent}{0pt}
\setlength{\parskip}{5pt}
\setlength{\emergencystretch}{3em}
\newcommand{\N}{\mathbb{N}}
\newcommand{\Z}{\mathbb{Z}}
\newcommand{\Q}{\mathbb{Q}}
\newcommand{\R}{\mathbb{R}}
\newcommand{\C}{\mathbb{C}}
\newcommand{\F}{\mathbb{F}}
\newcommand{\A}{\mathbb{A}}
\newcommand{\PP}{\mathbb{P}}
\newcommand{\E}{\mathbb{E}}
\newcommand{\eps}{\varepsilon}
\newcommand{\dd}{\,\mathrm{d}}
\begin{document}
\section*{UnsolvedMath ID: 20000408}
\textbf{Source problem code:} catalogue entry
\subsection{Definitions and mathematical statement}
The AIM item asks whether averages of Hardy--Littlewood constants can be computed in the function-field setting, followed by the word ``Yes''. It specifies no family, normalization or limit. The supplied report explicitly reconstructs one reading, without asserting it is the only intended one.

Fix a prime power $q$ and an integer tuple size $k\geq1$. Let $B_m=\{h\in\F_q[T]:\deg h<m\}$, including zero, with $N_m=q^m$. For an ordered pairwise-distinct tuple $\mathbf h=(h_1,\ldots,h_k)$, let $\nu_P(\mathbf h)$ be the number of its distinct residue classes modulo a monic irreducible $P$, and put $Q_P=q^{\deg P}$. The prime-tuple singular series in this normalization is
\[
 \mathfrak S(\mathbf h)=\prod_P\frac{1-\nu_P(\mathbf h)/Q_P}{(1-1/Q_P)^k}.
\]
A tuple is admissible when $\nu_P<Q_P$ for every $P$; otherwise its series is zero. Pairwise distinctness matters for the full product.

The report claims the quantitative degree-box mean
\[
 q^{-mk}\sum_{\substack{\mathbf h\in B_m^k\\h_i\ne h_j\ (i\ne j)}}\mathfrak S(\mathbf h)
 =1+O_{q,k}\bigl(1/\log(m+2)\bigr),
\]
and the same estimate when normalized instead by $(N_m)_k=N_m(N_m-1)\cdots(N_m-k+1)$. These are fixed-$q$, fixed-$k$, large-$m$ statements for the full series, not just a finite truncation.

For a finite set of primes, independent uniform residue tuples give an exact truncated mean of one by CRT. The associated local occupancy formula, for an integer moment $a\geq1$ and $Q=Q_P$, is
\[
 \E\left[\left(\frac{1-\nu_P/Q}{(1-1/Q)^k}\right)^a\right]
 =(1-1/Q)^{-ka}\sum_{j=1}^{\min(k,Q)}
 \frac{(Q)_j\left\{\begin{smallmatrix}k\\j\end{smallmatrix}\right\}}{Q^k}(1-j/Q)^a,
\]
where the braces are Stirling numbers of the second kind. This residue model permits repeated residues; exact CRT averaging does not by itself prove the distinct-tuple full-tail estimate.

The report also gives an $L^2$ profinite residue-product model with mean one. In that model the admissibility probability is the finite product
\[
 p_{\mathrm{adm}}(q,k)=\prod_{Q_P\leq k}
 \left(1-\frac{Q_P!\left\{\begin{smallmatrix}k\\Q_P\end{smallmatrix}\right\}}{Q_P^k}\right)>0.
\]
It claims the corresponding admissibility-conditioned limiting mean is $1/p_{\mathrm{adm}}$, not one. Finally it reports exponential convergence of the pair average, without supplying its exact rate or formula. The finite CRT identity, repeated-tuple treatment and full infinite-product tail are separate steps; detailed proofs and authoritative references are absent.
\subsection{Short English statement}
Compute averages of function-field Hardy--Littlewood constants after specifying a family. The report chooses fixed-field additive degree boxes and distinct ordered shifts, claiming full singular-series mean one with a quantitative error. It separates exact finite CRT moments from infinite-tail estimates and gives a different mean after conditioning on admissibility.
\subsection{Sources}
\begin{itemize}
\item Local registry: \url{data/unsolvedmath/problems.json}, permanent ID $20000408$, source code \texttt{catalogue entry}; original item, explicit-reconstruction qualification and supplied research summary. Statement status reconstructed and unverified.
\item AIM workshop ``Arithmetic statistics over finite fields and function fields'', Question 30, attributed as ``Rubenstein'' in the original: \url{https://aimath.org/pastworkshops/arithstatffieldproblems.pdf}. Canonical location \texttt{aim-algebraic-number-theory-notes.json}, \texttt{notes[59]}; original remarks and literature field empty. No PDF was opened here.
\item No full authoritative reference is supplied for the degree-box mean, occupancy/profinite construction or pair-average rate. The clean-statement field also contains the variant spelling ``Rubinstein''; no identity expansion is guessed.
\item Supplied machine-generated research attempt 1: partial result; problem solved at the run and in the registry, classification \texttt{SOLVED-IN-LITERATURE}; rigor metadata ``proved'', novelty confidence low. The non-peer-reviewed theorem package requires independent verification before citation.
\end{itemize}
\end{document}
